% ==========================================================
% 폰트 설정
% ==========================================================
%
% [기본] kotex가 제공하는 기본 한글 폰트를 그대로 사용합니다.
%        별도 설치 없이 Overleaf 및 로컬 TeX Live에서 바로 컴파일됩니다.
%        시각적으로 명조 계열 조판이 이루어지므로 초안용으로 충분합니다.
%
% [지침 준수] 가천대학교 지침은 "신명조 / 명조 / 휴먼명조 중 사용"을 요구합니다.
%            제출본에는 함초롬바탕 같은 실제 명조 폰트를 권장합니다.
%            아래 "옵션 A" 또는 "옵션 B" 섹션의 주석을 해제하세요.

% ----------------------------------------------------------
% 옵션 A) 함초롬바탕 (권장, 무료 배포, 지침 완전 준수)
% ----------------------------------------------------------
% 1) https://github.com/wkpark/HCR-LVT 또는
%    https://im.cm/hamchorom/ 등에서 함초롬바탕 .ttf 다운로드
% 2) thesis-gachon/fonts/ 폴더에 아래 2개 파일 배치:
%      - HANBatang.ttf
%      - HANBatangB.ttf
% 3) 아래 6줄 주석 해제
%
% \setmainhangulfont{HCR Batang}[
%   Path=./fonts/,
%   Extension=.ttf,
%   UprightFont=HANBatang,
%   BoldFont=HANBatangB
% ]

% ----------------------------------------------------------
% 옵션 B) 시스템에 설치된 명조 폰트 사용 (macOS/Windows)
% ----------------------------------------------------------
% macOS 예시:
% \setmainhangulfont{AppleMyungjo}
%
% Windows 예시:
% \setmainhangulfont{Batang}[BoldFont={Batang Bold}]

% ----------------------------------------------------------
% 영문 폰트 (옵션)
% ----------------------------------------------------------
% 지침은 국문 기준이지만, 표지/초록의 영문 제목 가독성을 위해
% 세리프 계열 영문 폰트를 지정하는 것을 권장합니다.
%
% \setmainfont{TeX Gyre Termes}   % Times 호환, TeX Live 기본 포함
